\documentclass[11pt,a4paper]{article}
\usepackage[T1]{fontenc}
\usepackage[utf8]{inputenc}
\usepackage{lmodern,amsmath,amssymb,amsthm,mathtools}
\usepackage[margin=25mm]{geometry}
\usepackage{booktabs,longtable,array,microtype,xcolor,enumitem,listings,hyperref}
\hypersetup{colorlinks=true,linkcolor=blue!45!black,urlcolor=blue!45!black,citecolor=blue!45!black,
 pdftitle={Smaller Diophantine forms of ZFC consistency},pdfauthor={Research contribution prepared for ProveIt}}
\setlength{\emergencystretch}{2em}
\setlength{\parskip}{3pt}
\setlist{itemsep=2pt,topsep=4pt}
\newtheorem{theorem}{Theorem}[section]
\newtheorem{lemma}[theorem]{Lemma}
\newtheorem{proposition}[theorem]{Proposition}
\newtheorem{corollary}[theorem]{Corollary}
\theoremstyle{remark}\newtheorem{remark}[theorem]{Remark}
\newcommand{\Z}{\mathbb Z}
\newcommand{\N}{\mathbb N}
\newcommand{\Pos}{\mathbb Z_{>0}}
\newcommand{\Con}{\operatorname{Con}}
\newcommand{\pc}{\operatorname{pc}}
\newcommand{\sgn}{\operatorname{sgn}}
\newcommand{\Positive}{\operatorname{Pos}}
\newcommand{\src}[1]{\nolinkurl{#1}}
\newcommand{\pin}{\texttt{3b201cdea32d2c816f5f3d20734051f25fd3ef14}}
\lstset{basicstyle=\ttfamily\footnotesize,breaklines=true,columns=fullflexible,keepspaces=true,frame=single}
\title{\textbf{Smaller Diophantine forms of ZFC consistency}\\[5pt]
\large Combining the blank-input compiler, ordinary Pell ranks,\\
and signed auxiliary coordinates}
\author{Research contribution prepared for the ProveIt repository}
\date{22 October 2026\\[3pt]\small Audited repository snapshot: \pin}
\begin{document}
\maketitle
\begin{abstract}
We combine three existing construction techniques in ProveIt and prove
additional domain reductions. With the fixed coefficients of package
\src{06-blank67}, an ordinary strong Pell equation gives a system of six
equations in fifteen positive unknowns, of maximum degree 28, evaluated
in 61 binary arithmetic operations. This improves the previous
16-unknown, seven-equation, degree-28, 63-operation system, and lowers the
degree of the previous 15-unknown, six-equation, 61-operation system.
The corresponding product polynomial has exact degree 99 and 69
operations; its sum-of-squares alternative has exact degree 56 and 78
operations. A quadratic lift uses 33 positive unknowns and 24 equations,
giving a single quartic in 33 positive unknowns.

A separate argument removes positivity from six auxiliary coordinates,
permits an odd strong multiplier, and makes the power scale a square.
Only seven coordinates still need a positive parametrization. This gives
explicit equations in 29 unrestricted integer unknowns, including a
121-operation polynomial of exact degree 188. Two alternative integer
polynomials use 119 and 120 operations at degrees 300 and 236.

For every displayed formulation, ZFC proves that absence of solutions
in its stated domain is equivalent to $\Con(\mathrm{ZFC})$. The proof
uses the repository's ordinary K215 calibration and native-window/Pell
compiler proofs; these dependencies are identified explicitly. The new
bridge is not a Lean or Rocq formalization. The four coefficients are
closed finite integer recipes, with $B=2^{5^{67}}$; their binary
expansions are not materialized. No global minimum, consistency proof,
or improvement to a varying-input universal operation bound is claimed.
\end{abstract}
\clearpage\tableofcontents\clearpage

\section{Results and the meaning of simplicity}
The user-specified objective has several different size measures:
number of unknowns, number of equations, degree, arithmetic circuit
length, and coefficient height. There is no agreed total ordering of
these measures. We give explicit improvements within the supplied
compiler family and preserve the alternatives needed to compare them
fairly. Our recommended certificate is the six-equation system below.
For a single equation in positive integers, the new degree-99 product
is a useful degree/operation compromise. For unrestricted integers, the
29-variable degree-188 product is the main new formulation.

All operation counts charge each binary $+$, $-$, and multiplication,
including multiplication by a fixed numeral. Fixed coefficients and
reuse of a computed register are free. A system's equality comparisons
are not arithmetic operations. Conversion to a single sum of squares
charges its subtractions, squares, and additions in full. These are
upper bounds for specified circuits, not lower bounds for all circuits.

\begin{center}\small
\begin{tabular}{lrrrr}
\toprule
New presentation & Unknowns & Equations & Degree & Operations\\
\midrule
Ordinary Pell system & $15$ positive & $6$ & $28$ max. & $61$\\
Product $P_{69}$ & $15$ positive & $1$ & $99$ & $69$\\
Sum of squares $S_{78}$ & $15$ positive & $1$ & $56$ & $78$\\
Quadratic lift & $33$ positive & $24$ & $2$ max. & ---\\
Quartic $Q_{33}$ & $33$ positive & $1$ & $4$ & ---\\
Integer system & $29$ integers & $6$ & $50$ max. & $113$\\
Integer product & $29$ integers & $1$ & $188$ & $121$\\
Integer sum of squares & $29$ integers & $1$ & $100$ & $130$\\
Integer normalized alternative & $29$ integers & $1$ & $300$ & $119$\\
Integer reduced alternative & $29$ integers & $1$ & $236$ & $120$\\
\bottomrule
\end{tabular}
\end{center}
All degrees in the table are exact, including the maximum residual
degrees. The artifacts give every arithmetic gate and every input port.

\begin{theorem}[Main consistency formulations]\label{thm:main}
Use the four fixed coefficients in Appendix~\ref{app:coefficients},
bound to the exact K215 table in the archive. Let $\mathcal S$ be the
system in Section~\ref{sec:formulas}, let $P_{69}$ and $S_{78}$ be its
displayed single-polynomial forms, and let $Q_{33}$ and $\widehat P_{121}$
be the quartic and unrestricted-integer formulations below. Then ZFC
proves each of the equivalences
\begin{align*}
 \exists\mathbf x\in\Pos^{15}\;\mathcal S(\mathbf x)
 &\ \longleftrightarrow\ \exists\mathbf x\in\Pos^{15}\;P_{69}(\mathbf x)=0\\
 &\ \longleftrightarrow\ \exists\mathbf x\in\Pos^{15}\;S_{78}(\mathbf x)=0\\
 &\ \longleftrightarrow\ \exists\mathbf v\in\Pos^{33}\;Q_{33}(\mathbf v)=0\\
 &\ \longleftrightarrow\ \exists\mathbf t\in\Z^{29}\;\widehat P_{121}(\mathbf t)=0\\
 &\ \longleftrightarrow\ \exists p\;\operatorname{Proof}_{\mathrm{ZFC}}(p,\bot).
\end{align*}
Consequently, the negation of each solution predicate is provably
equivalent in ZFC to $\Con(\mathrm{ZFC})$.
\end{theorem}
Here $\operatorname{Proof}_{\mathrm{ZFC}}$ is a standard effective
finite proof predicate, or the equivalent explicit axiom presentation
used by the repository's calibration. Standard finite proof
translations between presentations are formalizable in ZFC.

\section{What the eight earlier packages contribute}
The comparison is pinned to \pin. The eight packages themselves share
the earlier snapshot \src{cdd582bec13f908df3711fa9e05d1ddac368a5db}.
The repository reconciliation records their differing coefficients,
proof status, and metrics~\cite{reconciliation}. The following table
selects the constructions relevant to the present improvement; it is
not a claim to reproduce every prior artifact.

\begin{center}\small
\begin{tabular}{p{29mm}p{71mm}p{34mm}}
\toprule
Package & Relevant contribution & Selected prior size\\
\midrule
01-explicit82 & Direct fixed-input specialization with retained witnesses & $18$ vars., $82$ ops.\\
02-system63 & Ordinary strong-root equation and reduced auxiliary coefficient & $16/7/28/63$\\
03-system64 & Equivalent fixed-marker system; strong-root elimination & $16/7/28/64$\\
04-poly73 & Affine marker guard and two-sided normalization & $16$ vars., $73$ ops.\\
05-fixed15-16 & Witness restoration and strong-root elimination alternatives & $15$ vars., $80$ or $82$ ops.\\
06-blank67 & Raw K215, no input interface, numeric coefficient recipe & $15/6/60/61$; $P_{67}$\\
07-fixed70 & Fixed input with altered Pell scale & $17/8/28/61$\\
08-start-only67 & Independent input-free arithmetic; quadratic register elimination & $34$ vars., $25$ quadratics\\
\bottomrule
\end{tabular}
\end{center}
In four-number entries the order is unknowns/equations/degree/operations.
Packages 06 and 08 independently have the same 15-variable,
67-operation, degree-155 polynomial skeleton, but different coefficient
integers. We retain \emph{only package 06's coefficient recipe}; borrowing
the ordinary-root algebra of 02 or the register eliminations of 08 does
not authorize substituting their different coefficients.

The new system improves the $16/7/28/63$ system on unknowns, equations,
and operations, while retaining degree 28. It improves the $15/6/60/61$
system on degree without increasing any other listed size. The new
quartic reduces the earlier $34$ positive variables and $25$ quadratic
relations to $33$ and $24$. The degree-99 product costs two operations
more than $P_{67}$; we do not describe it as beating the operation record.
The earlier degree-123 product costs 68 operations. The new degree-56
sum of squares costs 78 operations, versus 83 for package 02's
16-variable degree-56 form and 79 for package 06's degree-92 form.

Fifteen is not the least known possible number of positive/natural
unknowns in general Diophantine representation. Jones's classical
nine-unknown construction has a very different degree and description
cost~\cite{jones}. Similarly, 29 is a size bound for these fully specified
integer equations, not a universal minimum over all integer
representations; see also~\cite{bayer}. The present gain is a concrete
combination and proof simplification relative to the supplied artifacts.

\section{The six equations and a single polynomial}\label{sec:formulas}
Let $\beta,K,M_C,M_F^*$ be the fixed coefficients of
Appendix~\ref{app:coefficients}. In particular, $\beta=B-1$, and
$M_F^*$ is the \emph{shifted} source coefficient $M_F+B-1$.
Supply exactly the fifteen positive integers
\[
 (J,F,\alpha,z,f,h,i,T,s,w,\tau,\eta,\zeta,y,\gamma).
\]
Define polynomial abbreviations
\begin{align}
 q&=\beta J+1,& C&=q-F-\alpha,& X&=wq,&Y&=sq^3,\notag\\
 E&=XY,&k&=\eta+\zeta,&c&=kY+\eta,&a&=E+Y,\notag\\
 H&=4a+3,&\Delta&=a^2+H=(a+2)^2-1,&
 D&=X+ac+\gamma H.\label{eq:abbrev}
\end{align}
Further put
\begin{align}
 G&=(q-1)(q-F)+\alpha,\notag\\
 R&=G(q^2-1)+(M_C+qM_F^*)J,\notag\\
 L&=EkY,&V&=c(Tf-1)-Rf^2,&I&=ic^2,&U&=\Delta(f^2-1).
 \label{eq:packing}
\end{align}
These are substitutions, not additional unknowns or constraints.
All exponents of unknowns are fixed nonnegative integers.
Set
\begin{align}
 N_0&=\tau^2-L(L+k),&N_1&=D^2-\Delta c^2,\notag\\
 N_a&=U(V^2-y^2)+y^2,&N_k&=k-hE-R,\notag\\
 N_t&=(K+w)C+q-F-z(q-1),&N_s&=1+I^2-U.
 \label{eq:factors}
\end{align}
Our recommended system is
\[
 \boxed{N_0=N_1=N_a=N_k=N_t=1,\qquad I^2=U.}
 \tag{S}\label{eq:system}
\]
The degree-99 single equation is
\[
 \boxed{P_{69}=N_0N_1N_aN_kN_tN_s-1=0.}
 \tag{P}\label{eq:product}
\]
The degree-56 alternative is
\[
 \boxed{S_{78}=(N_0-1)^2+(N_1-1)^2+(N_a-1)^2
 +(N_k-1)^2+(N_t-1)^2+(I^2-U)^2=0.}
 \tag{SOS}\label{eq:sos}
\]
This is an ordinary polynomial problem: the coefficient recipe has no
variable program parameter, hidden exponential relation, extraction
function, or variable-length equation list.

\section{Pell preliminaries and the ordinary rank lemma}
For an integer $A\ge2$, write
\[
 \chi_A(n)+\psi_A(n)\sqrt{A^2-1}
 =(A+\sqrt{A^2-1})^n\qquad(n\in\N).
\]
Both sequences are integer recurrences with characteristic equation
$t^2-2At+1=0$. Positive solutions of
$d^2-(A^2-1)c^2=1$ have $(d,c)=(\chi_A(p),\psi_A(p))$.
We use the addition and composition identities, the growth bounds
\[
 (2A-1)^{n-1}\le\psi_A(n)\le(2A)^{n-1},
\]
for $n\ge1$, and strong divisibility
$\gcd(\psi_A(u),\psi_A(v))=\psi_A(\gcd(u,v))$.
Strong divisibility follows from the addition formula,
$\gcd(\chi_A(n),\psi_A(n))=1$, and the Euclidean algorithm on indices.
These assertions concern integer-coefficient units; no assertion about
the full ring of algebraic integers is needed.

\begin{lemma}[Ordinary strong rank]\label{lem:ordinary}
Let $\Delta=A^2-1$, $c=\psi_A(p)$, and $c>A\Delta^2$.
If nonzero integers $i,f$ satisfy
\[
 (ic^2)^2=\Delta(f^2-1),
\]
then there is $m>0$ with
\[
 |f|=\chi_A(m),\qquad p\mid m,\qquad c\mid m,
 \qquad |i|c^2=\Delta\psi_A(m).
\]
In particular $c^2\mid f^2-1$.
\end{lemma}
\begin{proof}
We include the inherited rank argument~\cite{rank} because it is the
essential justification for the change of strong equation.
Write $\Delta=d_0s_0^2$, with $d_0>1$ squarefree, and let
$F_0+y_0\sqrt{d_0}$ be the least integer-coefficient norm-one unit
greater than 1. Division by successive powers of this least unit proves that
all positive integer solutions in this field are its powers. Thus for
some $e\ge1$, with $S=\psi_{F_0}(e)$,
\[
 A=\chi_{F_0}(e),\quad s_0=y_0S,\quad
 \Delta=(F_0^2-1)S^2,\quad\psi_{F_0}(ep)=Sc.
\]
Set $\rho=|i|c^2$. The equality $\rho^2=\Delta(f^2-1)$ first gives
$s_0\mid\rho$, then $d_0\mid\rho/s_0$ by squarefreeness. Hence
$|f|^2-d_0(\rho/(d_0s_0))^2=1$, and for some $n>0$
\[
 |f|=\chi_{F_0}(n),\qquad
 \rho=(F_0^2-1)S\psi_{F_0}(n).
\]
Let $T_0=(F_0^2-1)S\le\Delta$ and $g=\gcd(n,ep)$.
Strong divisibility, together with $c^2\mid T_0\psi_{F_0}(n)$, gives
\[
 c\le\gcd(c^2,\Delta c)
 \le T_0\psi_{F_0}(g),
 \qquad \psi_{F_0}(g)\ge c/\Delta.
\]
If $g<ep$, then $2g\le ep$. Since
$\psi_{F_0}(2g)>2\psi_{F_0}(g)^2$ and $S<A$, we obtain
$\psi_{F_0}(g)^2<Ac/2$. This contradicts
$c^2/\Delta^2>Ac$. Therefore $ep\mid n$; writing $n=em$
gives $p\mid m$, $|f|=\chi_A(m)$, and
$\rho=\Delta\psi_A(m)$.

Write $m=pb$ and $d=\chi_A(p)$. Expanding
$(d+c\sqrt\Delta)^b$ modulo $c^2$ gives
$\psi_A(pb)/c\equiv b d^{b-1}\pmod c$. Thus
$c\mid\Delta b$. Also the binomial expansion modulo $\Delta$
gives $c\equiv pA^{p-1}\pmod\Delta$, so
$\gcd(c,\Delta)=\gcd(p,\Delta)$ divides $p$.
Consequently $c\mid pb=m$. Finally
$f^2-1=\Delta\psi_A(m)^2=|i|c^2\psi_A(m)$.
\end{proof}

\begin{lemma}[Two unit exclusions]\label{lem:units}
If $v>1$ is an integer, $x^2-v(v+1)t^2=-1$ has no integer solution.
If $\Delta=A^2-1$, neither $D^2-\Delta c^2$ nor
$1+(ic^2)^2-\Delta(f^2-1)$ can equal $-1$.
\end{lemma}
\begin{proof}
For the first assertion take a solution with $t>0$ minimal and $x>0$.
Its equation implies $vt<x<(v+\tfrac12)t$. Multiplication by
$2v+1-2\sqrt{v(v+1)}$ gives another integer solution whose coefficient
$t'=(2v+1)t-2x$ satisfies $0<t'<t$, a contradiction.
For the others, $\Delta$ is $0$ or $3$ modulo 4. The first expression
cannot be $3$ modulo 4. For the second, note that
$\Delta(f^2-1)$ is $0$ or $1$ modulo 4; subtracting it from a square
cannot give $-2$ modulo 4.
\end{proof}

\section{Soundness, with the proof order made explicit}\label{sec:soundness}
We give a proof that also covers the signed-domain extension. Initially
require $J,F,\alpha,z,s,\eta,\zeta>0$, $w\ge0$, and $i\ne0$.
Allow $f,h,T,\tau,y,\gamma$ to be arbitrary integers. For now assume
that either \eqref{eq:system} or \eqref{eq:product} holds.

\subsection{Unit factors and the zero-scale boundary}
At a product zero every factor is an integer unit. The mod-4 arguments
give $N_1=N_s=1$. Once $N_s=1$, $U=I^2$ is a square and
$N_a$ is congruent either to $y^2$ or to $V^2$ modulo 4;
therefore $N_a=1$. For $X>0$, the first norm is
\[
 N_0=\tau^2-v(v+1)k^2,\qquad v=XY^2>1,
\]
so $N_0=1$ by Lemma~\ref{lem:units}. If $X=0$, it is simply
$\tau^2$ and again cannot be $-1$.
Thus $N_k=\varepsilon$, $N_t=\nu$, where
$\varepsilon,\nu\in\{1,-1\}$ and $\varepsilon\nu=1$.

If $w=0$, then $X=0$, $a=Y>0$, $D=ac+\gamma H$, and
$\Delta=a^2+H$. The main norm reduces modulo $H$ to
\[
 D^2-\Delta c^2\equiv0\pmod H,
\]
contradicting its value 1 because $H=4Y+3>1$.
Hence $w>0$. This short modular argument is the reason a square
parametrization of $w$ does not require a separate nonzero constraint.

\subsection{Bounds before decoding any mask}
Since $z,F\ge1$,
\[
 (q-F)-z(q-1)=1-F-(z-1)(q-1)\le0.
\]
With $K+w\ge2$, $N_t\in\{1,-1\}$ therefore forces $C\ge0$.
It does \emph{not} yet force $C>0$. It follows that
$q-F=C+\alpha>0$ and $1\le F\le q-1$.
Using the authentic masks $0<M_C,M_F<B-1$ and
$M_F^*=M_F+B-1$, one gets
\begin{align}
 q&\le G=q\alpha+(q-1)C\le q^2-q,\notag\\
 0&<(M_C+qM_F^*)J<2q^2-q-1,\notag\\
 q(q^2-1)&<R<q^4-q^3+q^2-1<q^4-2,\notag\\
 q&\ge B\ge16,\quad X\ge q,\quad Y\ge q^3,
 \quad E\ge q^4>R+2.\label{eq:bounds}
\end{align}
In particular $R>3q+1$ and $a>R+2$. None of these conclusions
uses a dyadic value of $q$, an input equation, or $C>0$.

\subsection{First and main Pell indices}
Put $A=a+2$ and $P_0=2XY^2+1$. The least integer-coefficient
norm-one unit greater than 1 for $v(v+1)$ is $2v+1+2\sqrt{v(v+1)}$:
coefficient 1 cannot solve the norm, and coefficient 2 does.
Consequently
\[
 k=2\psi_{P_0}(n),\quad |\tau|=\chi_{P_0}(n),\qquad
 c=\psi_A(p),\quad |D|=\chi_A(p)
\]
for $n,p\ge1$. Since $P_0\equiv1\pmod E$,
\begin{equation}
 2n=R+\varepsilon+v_1E,\qquad v_1\ge0,
 \qquad n\ge(R-1)/2\ge24.
 \label{eq:first-index}
\end{equation}
The nonnegativity of $v_1$ uses $0<R+\varepsilon<E$.
Since $P_0>A$ and $c>k$, monotonicity of the Pell sequences gives
$p>n$. Their growth then gives
\[
 c>A\Delta^2,\qquad c>2p.
\]
Even when $h$ was allowed to have either sign,
$k=R+\varepsilon+hE>0$ and \eqref{eq:bounds} imply $h\ge0$.
Thus $c>kY\ge(R-1)Y>2R$.
In fact $h=0$ is impossible: it would give $k=R\pm1$;
$n=1$ gives $k=2$, whereas $n\ge2$ gives
$k\ge4P_0>R+1$. We may therefore recover $h>0$, although
the weaker bound already suffices.

\subsection{Signed auxiliary congruences recover the exact index}
Lemma~\ref{lem:ordinary} now applies, before any sign restriction on
$V$ or $T$ is used. Write $f_0=|f|$ and
$\rho=|i|c^2>1$. We obtain
\[
 f_0=\chi_A(m),\quad c\mid m,\quad m\ge c>2p,
 \quad\rho=\Delta\psi_A(m),\quad c^2\mid f^2-1.
\]
The auxiliary norm is
\[
 \rho^2V^2-(\rho^2-1)y^2=1.
\]
It forces $V\ne0$ and $y\ne0$. Pell classification gives
$\rho|V|=\chi_\rho(\ell)$ and $|y|=\psi_\rho(\ell)$.
The index $\ell$ is odd, since an even-index $\chi_\rho$ is
$\pm1$ modulo $\rho$. Write $\ell=2u+1$ and
$\sigma=\sgn(V)$. The integer polynomials $Q_u$ determined by
\[
 \chi_Z(2u+1)=ZQ_u(Z^2)
\]
satisfy
\[
 Q_u(0)=(-1)^u(2u+1),\qquad
 Q_u(1-A^2)=(-1)^u\psi_A(2u+1).
\]
For example $Q_0=1$, $Q_1=4t-3$, and
$Q_{u+2}=(4t-2)Q_{u+1}-Q_u$; the identities follow by recurrence.

The definition of $V$ supplies, without any auxiliary positivity proof,
\[
 V\equiv-c\pmod{f_0},\qquad V\equiv-R\pmod c.
\]
The second congruence uses the rank conclusion $c^2\mid f^2-1$;
it must not be obtained by illicit cancellation of $\Delta$ in the
ordinary strong equation. Since $\rho^2\equiv1-A^2\pmod{f_0}$
and $c\mid\rho$, we have
\[
 \sigma(-1)^u\psi_A(\ell)\equiv-c\pmod{f_0},\qquad
 \sigma(-1)^u\ell\equiv-R\pmod c.
\]
Squaring the first congruence and using
$\chi_A(2j)=1+2\Delta\psi_A(j)^2$ yields
\[
 \chi_A(2\ell)\equiv\chi_A(2p)\pmod{\chi_A(m)}.
\]
The Pell step-down lemma~\cite{stepdown}, applicable to $0<2p<m$,
gives $\ell\equiv\pm p\pmod m$. A short proof reduces indices
using $\chi_A(j+2m)\equiv-\chi_A(j)\pmod{\chi_A(m)}$ and
evenness, and distinguishes representatives by strict monotonicity
and $2\chi_A(m-1)<\chi_A(m)$.
As $c\mid m$, the other congruence implies $R\equiv\pm p\pmod c$.
The bounds $0<R,p<c/2$ force
\[
 \boxed{p=R.}
\]
This is the key proof simplification: it eliminates the need to restore
positive quotient coordinates before identifying the index.

\subsection{No wrap, and the two remaining signs}
If $v_1\ge1$ in \eqref{eq:first-index}, then
$2n\ge R-1+E>2R$, contradicting $n<p=R$.
Hence $2n=R+\varepsilon$. If $\varepsilon=-1$, then $p=2n+1$.
With $Q=\chi_A(2)>P_0$, duplication and monotonicity imply
\[
 \psi_A(2n)=2A\psi_Q(n)\ge2A\psi_{P_0}(n)=Ak.
\]
This contradicts $c=\psi_A(2n+1)<k(Y+1)$, since $A>Y+1$.
Therefore $\varepsilon=\nu=1$, $R=2n-1$ is odd, and
$N_t=1$ now gives $C>0$. Every product zero satisfies the six
equations individually. The converse is immediate, proving their
equality of zero sets on the positive domain.

\subsection{The sign of the main root and power recovery}
Put $r_0=(R-1)/2\ge2$. The inherited lower Pell-ratio estimate
uses only $c=\psi_A(R)$ and $k=2\psi_{P_0}(r_0+1)$:
\[
 \frac{c}{k/2}>
 \frac{(X+1)^{2r_0}}{X^{r_0}}>X^{r_0}.
\]
For completeness, the stronger first comparison follows by bounding
the numerator below by $(2A-1)^{2r_0}$ and the denominator above by
$(2P_0)^{r_0}$; the extra factor is greater than one because
$6XY^2>a$. Together with $c/k<Y+1$, it gives
\[
 Y>X^{r_0}/3,\qquad a>X^{r_0+1}/3.
\]
As $X\ge16$ and $r_0\ge2$, $(X/4)^{r_0}>6$, and hence
\begin{equation}
 0<X2^R=2X4^{r_0}<a<H-1.\label{eq:strongratio}
\end{equation}
The recurrence gives $\chi_A(R)-a\psi_A(R)\equiv2^R\pmod H$.
Also
$(\chi_A(R)-ac)(\chi_A(R)+ac)=1+Hc^2$.
If $D=-\chi_A(R)$, the defining relation
$D-ac=X+\gamma H$ would imply $X2^R\equiv-1\pmod H$,
contradicting \eqref{eq:strongratio}. Thus $D>0$ even for signed
$\gamma$. We conclude $X\equiv2^R\pmod H$; both numbers lie
strictly between zero and $a$, so
\[
 \boxed{X=2^R.}
\]
Since $q\mid X$, $q=2^t$. The relation $q\equiv1\pmod{2^d-1}$
gives $d\mid t$ by the order-of-two argument; write $q=B^N$.
The bound $R>3q$ implies $q^3\mid X$. Thus the post-index
ratio and binomial projection calculation in Sections 5--6 of
\src{pell_kernel_half_binomial42.md} applies to the recovered
positive $X,Y,c,k$, with $p=R$, $2n=R+1$, and the strict ratio
$kY<c<k(Y+1)$. This calculation does not use the signs of the
remaining auxiliary coordinates. We reuse that part of the proof,
not its full positive-witness theorem on a signed tuple.

\subsection{The retained computation compiler}
At this point the inherited binomial calculation gives
$Y=\lfloor (X+1)^{R-1}/X^{(R-1)/2}\rfloor/2$ and
$\nu_2(Y)=\pc(R)-2$. The paid factor $q^3\mid Y$ supplies the
population threshold used by the inverse mask theorem. The resulting
relations are
\[
 (C-1)\mathbin{\&}(M_CJ+1)=0,\qquad
 F\mathbin{\&}(M_FJ-1)=0.
\]
Here $\&$ is used only in this proof, not in the polynomial.
The carry-free convolution and anchor argument of the retained
compiler extract a cyclic word of local windows with Start at the
origin and the distinguished impossible selector excluded. Start's
top/center mismatch rules out zero temporal shift. The blank-window
lemma extracts a finite halting computation of the literal K215 table.
The exact compiler proof dependencies and the local language are
included in the archive; Appendix~\ref{app:coefficients} fixes all
coefficients. This inherited decoding stage is unchanged by the
ordinary-root or sign substitutions.

\section{Completeness and the ZFC equivalence}\label{sec:complete}
Suppose K215 has a finite blank halting run. The retained compiler pads
it to a helical rectangle with spatial width and temporal height both
large powers of five. It constructs the outer coordinates and the
first and main Pell witnesses for the actual packed index $R$.
The dummy-bit adjustment arranges the required temporal congruence
at that same $R$; it is not an earlier provisional index. In particular
\[
 q=B^N,\quad N\text{ a power of five},\quad X=2^R,
 \quad R\equiv3\pmod4,
\]
and all originally positive coordinates can be chosen positive.
The already-proved normalized construction has
$f^2-1=\Delta(i_{\rm norm}c^2)^2$.
Set $i_{\rm ord}=\Delta i_{\rm norm}$. Then
\[
 (i_{\rm ord}c^2)^2=\Delta(f^2-1),
\]
and the auxiliary coefficient and all other equations are unchanged
on the solution. This proves completeness of \eqref{eq:system}.
It proves equivalence of \emph{existence} predicates, rather than
identity of tuples with the same strong multiplier.

The existing K215 calibration states
\[
 \operatorname{Halt}(\mathrm{K215},\mathrm{blank})
 \ \longleftrightarrow\ \exists p\,
 \operatorname{Proof}_{\mathrm{ZFC}}(p,\bot).
\]
Its ordinary proof, the finite syntax and computation encodings, the
Pell and coefficient arguments, and the present transformations are
formalizable in ZFC. Composing these arguments proves
Theorem~\ref{thm:main}. This reasoning does not infer ZFC-provability
from the metatheoretic soundness of Lean or Rocq.

\section{Twenty-nine unrestricted integer unknowns}\label{sec:integer}
The signed soundness proof above already established more than
positive-domain equivalence. It permits arbitrary integer values for
\[
 f,h,T,\tau,y,\gamma,
\]
provided $i\ne0$, $w\ge0$, and
$J,F,\alpha,z,s,\eta,\zeta>0$. We next obtain convenient polynomial
parametrizations of $i$ and $w$ without losing completeness.

\subsection{An odd strong multiplier is sufficient}
In a canonical halting witness, $A=a+2$ is even and $R$ is odd;
therefore $c=\psi_A(R)$ is odd. Choose the auxiliary index
\[
 m=cR
\]
instead of the older $2cR$. Expanding
$(\chi_A(R)+c\sqrt\Delta)^c$ shows $c^2\mid\psi_A(cR)$.
Since $m$ is odd and $\psi_A(m)\equiv m\pmod2$, the integer
\[
 i_{\rm norm}=\psi_A(cR)/c^2
\]
is positive and odd. Also $\Delta$ is odd, so
$i_{\rm ord}=\Delta i_{\rm norm}$ is positive and odd.
Set $f=\chi_A(m)$ and $\rho=\Delta\psi_A(m)$, and take
\[
 V=\chi_\rho(R)/\rho,\qquad y=\psi_\rho(R).
\]
The quotient is an integer because $R$ is odd. As $R\equiv3\pmod4$,
the identities for $Q_u$ give
$V\equiv-c\pmod f$ and $V\equiv-R\pmod c$.
The coprimality $\gcd(c,f)=1$ and $f^2\equiv1\pmod c$ make
\[
 T=\frac{V+c+Rf^2}{cf}
\]
a positive integer. Thus the whole auxiliary block, not merely its
strong equation, has a witness with odd $i$.
We may replace $i$ everywhere by $2r+1$, with $r\in\Z$.
It is never zero, and covers the canonical positive odd choices.

\subsection{The power scale can be a square}
The fixed $d=5^{67}$ and canonical $N$ are odd, and $R$ is odd.
Moreover $R>dN$. Therefore
\[
 w=X/q=2^{R-dN}
   =\left(2^{(R-dN)/2}\right)^2.
\]
We may replace $w$ by $u^2$, $u\in\Z$. Soundness at $u=0$
was already excluded by the main norm modulo $H$; no extra equation
or inequality is needed.

\subsection{The last seven positive coordinates}
Define
\[
 \Positive(a,b,c)=1+a^2+b^2+c(c+1)\qquad(a,b,c\in\Z).
\]
It is always positive. Conversely every $n>0$ is represented:
$4(n-1)+1$ is not of the excluded three-square form
$4^e(8t+7)$, and is $1$ modulo 4. Legendre's three-square theorem
therefore gives exactly one odd square and two even squares. Writing
them as $(2c+1)^2,(2a)^2,(2b)^2$ yields the formula.

For each $v\in\{J,F,\alpha,z,s,\eta,\zeta\}$ use a separate
triple $(v_1,v_2,v_3)$ and substitute
\begin{equation}
 v=\Positive(v_1,v_2,v_3),\qquad i=2r+1,\qquad w=u^2.
 \label{eq:integer-substitution}
\end{equation}
Keep $f,h,T,\tau,y,\gamma$ as unrestricted integers.
There are exactly $7\cdot3+2+6=29$ unknowns.
Applying these substitutions to $P_{69}$ defines the explicit
polynomial $\widehat P_{121}$. Its lack of an integer zero is
provably equivalent in ZFC to $\Con(\mathrm{ZFC})$.

The substitutions cost $7\cdot7+2+1=52$ operations:
each positive parametrization costs three multiplications and four
additions, the odd affine form costs two, and the square costs one.
Thus $P_{69}$ becomes a 121-operation polynomial; the system becomes
113 operations; the sum of squares becomes 130 operations.
The same signed proof applies to the normalized $P_{67}$ and
reduced $P_{68}$ of package 06, giving 119- and 120-operation
29-integer alternatives. For normalized rank one may use the shorter
direct argument $\psi_A(m)=|i|c^2\Rightarrow pc\mid m$;
the ordinary relaxed-rank lemma is then unnecessary.

\section{A quartic in thirty-three positive unknowns}\label{sec:quartic}
For a low-degree presentation, retain the fifteen original positive
unknowns and introduce the following eighteen positive registers:
\[
 (n_3,E,p_0,c,a,X,D,c_2,G,\ell_2,R,f_2,v_0,V,g,I,b_0,\Delta).
\]
In the equations below $q=\beta J+1$, $k=\eta+\zeta$, and
$H=4a+3$ are linear abbreviations, not additional unknowns.
Take the following 24 quadratic equations:
\begin{gather}
\begin{aligned}
 \ell_2&=q^2-1,& n_3&=(\ell_2+1)q,\\
 X&=wq,&a&=E+n_3s,\\
 E&=X(a-E),&c&=\eta+(a-E)k,\\
 p_0&=E(c-\eta)+k,&\tau^2&=p_0(p_0-k)+1,\\
 D&=X+ac+\gamma H,&D^2&=\Delta c_2+1,\\
 \Delta&=a^2+4a+3,&c_2&=c^2,\\
 k&=Eh+R+1,&G&=(q-1)(q-F)+\alpha,
\end{aligned}\notag\\
 R=G\ell_2+(M_C+qM_F^*)J,\notag\\
 (K+w)(q-F-\alpha)+q-F-z(q-1)=1,\notag\\
\begin{aligned}
 f_2&=f^2,&v_0&=Tf-1,\\
 V&=cv_0-Rf_2,&g&=y^2-V^2,\\
 I&=ic_2,&b_0&=I^2,\\
 b_0&=\Delta(f_2-1),&y^2-b_0g&=1.
\end{aligned}\label{eq:quadratic}
\end{gather}
Every displayed polynomial has degree at most two after inlining the
three linear abbreviations. Let $r_1,\ldots,r_{24}$ be the left sides
minus the right sides. Then
\[
 \boxed{Q_{33}=\sum_{j=1}^{24}r_j^2=0}
\]
has exact degree four; the coefficient of $\tau^4$ is 1.

To verify equivalence, the first equations force
$n_3=q^3$, $a-E=sq^3=Y$, $E=XY$, and $c=kY+\eta$.
The register $p_0=EkY+k=L+k$ gives the original first norm.
All other substitutions recover \eqref{eq:system} directly.
For the converse, a positive system solution has positive values for
all these registers. The only nontrivial ones are $v_0,V,g$.
Lemma~\ref{lem:ordinary} and $c>A\Delta^2$ give $f>2c$ and
$f^2>\Delta+c$. Since $R<a<\Delta$,
\[
 U-(Rf^2+c)=(\Delta-R)f^2-\Delta-c>0.
\]
If $|V|>1$, the auxiliary norm implies $|V|\ge2U-1$;
but $T>0$ gives $V>-Rf^2-c>-U$. The values $V=\pm1$
are excluded by $V\equiv-c\pmod f$ and $2<c<f/2$.
Thus $V>1$, $v_0=Tf-1>0$, and the auxiliary norm gives
$g=y^2-V^2>0$.
This proves a positive inverse lift.

Relative to the 34-variable lift of package 08, the removed register
was the scaled quantity $\Delta c^2$. Its two defining norm equations
are replaced by $D^2=\Delta c_2+1$, and the strong register becomes
$I=ic_2$. These are algebraic changes; package 06's coefficients
remain the ones used here. Substituting $v=n_v+1$ for each of the
33 positive unknowns gives a degree-four equation over $\N^{33}$.
It is not legitimate simply to read $Q_{33}=0$ over unrestricted
integers without imposing its domain.

\section{Exact degrees and reproducible operation counts}\label{sec:counts}
The ordinary system replaces four gates of the normalized system
by four gates of the same cost:
\begin{center}\small
\begin{tabular}{ll}
\toprule
Old gates & New gates\\
\midrule
$i(\Delta c^2)$ & $f^2-1$\\
$(i\Delta c^2)^2$ & $\Delta(f^2-1)$\\
$\Delta f^2$ & $ic^2$\\
$\Delta f^2-(i\Delta c^2)^2$ & $(ic^2)^2$\\
\bottomrule
\end{tabular}
\end{center}
Each column costs three multiplications and one subtraction; the
new last two registers are compared as $I^2=U$. The result remains
$61=34M+27A$. Forming $N_s$ costs two operations, and the product
and subtraction of 1 cost six, giving $69=39M+30A$.
Six residual subtractions, six squares, and five additions give
$78=40M+38A$ for the sum of squares.

For the degree calculation put $Q=\beta J$, $k_0=\eta+\zeta$,
and $C_1=Q-F-\alpha$. The exact factor degrees of
$(N_0,N_1,N_a,N_k,N_t,N_s)$ are
\[
 (22,18,28,7,2,22).
\]
The only substantial top-degree cancellation occurs in $N_1$:
\[
 N_1=2ac(X+\gamma H)+(X+\gamma H)^2-Hc^2.
\]
Its leading term is $8\gamma a_{\rm top}^2c_{\rm top}$.
The product's full leading homogeneous form is
\begin{equation}
 8Q^{55}h\gamma i^2k_0^9w^7s^{16}T^2f^4
       \bigl(wC_1-zQ\bigr).
 \label{eq:leader}
\end{equation}
In particular
$J^{56}h\gamma i^2\eta^9w^7s^{16}T^2f^4z$
has coefficient $-8\beta^{56}\ne0$.
This proves degree 99 for the actual fixed coefficients, not just
for a synthetic coefficient example. The auxiliary residual is
uniquely highest, of degree 28, so the sum of squares has degree 56.

For the 29-integer substitution give weight 2 to
$J,F,\alpha,z,s,\eta,\zeta,w$ and weight 1 to the other old
coordinates. The factor weights become
\[
 (44,35,50,13,4,42),
\]
whose sum is 188. The leading form \eqref{eq:leader} remains the
weighted leader. Substitute each disjoint positive quadratic form,
$w=u^2$, and the leading term $2r$ of $i=2r+1$.
The result is nonzero: selecting the first square from every positive
form and the $-zQ$ term gives coefficient $-32\beta^{56}$.
The integer system has maximum degree 50 and its SOS degree 100.
The analogous weighted leaders for normalized $P_{67}$ and $P_{68}$
give degrees 300 and 236, respectively. The full calculations are in
\src{proofs/integer29_degrees.md}.

\section{Verification, provenance, and limits of the claim}
The archive is designed to run without network access. The arithmetic
generators and auditors use the Python standard library. They verify
the exact gate ledgers, absence of dead gates and free ports,
agreement of factored and expanded quadratic relations, and the
integer substitutions by exact arithmetic. Polynomial evaluations on
affine lines supplement the uniform leading-form proofs; they do not
replace them. Finite Pell checks test the identities and congruences
used in the domain reductions, and the inherited local-window checks
exercise actual marked halting control computations and invalid
truncated runs. These tests do not exhibit a halting K215 run.

Every copied source has its repository path and SHA-256 in
\src{SOURCE_MANIFEST.json}. The K215 table hash is
\begin{center}\footnotesize
\texttt{fdbcef0825886035c3091c8537a69535bc96f9e265f0ae45863543d40ae70090}.
\end{center}
The pinned endpoint README records Lean/Rocq checks for that same
literal table. Those builds were not rerun for this report. No new
theorem in either assistant is claimed. In particular the following
evidence levels should remain distinct:
\begin{enumerate}
\item the existing external proof-assistant theorems for K215;
\item the ordinary K215 calibration and compiler proofs, formalizable
      in ZFC and used as mathematical dependencies here;
\item the new ordinary proofs in this report;
\item finite source, arithmetic, and layout audits.
\end{enumerate}
An explicitly encoded derivation in ZFC's own proof syntax is not
provided. Nevertheless, using ordinary mathematical proofs carried
out within the resources of ZFC is the usual basis for the assertion
``ZFC proves the equivalence.'' No assumption of $\Con(\mathrm{ZFC})$
is used in proving the equivalence itself.

The four coefficients are astronomically large. Their fixed recipes
define actual integers, but the coefficient-height cost is not charged
in the circuit metric. No binary expansion of $B=2^{5^{67}}$, no
enumeration of all approximately $1.42\cdot10^{21}$ windows, and no
complete Diophantine witness is generated. A recipe of modest file
size must not be confused with small explicit coefficient values.

If ZFC is consistent, it cannot prove the asserted absence of
solutions, by the equivalence and the second incompleteness theorem.
Mere consistency is insufficient to claim that ZFC also cannot prove
the existence of a solution; such a two-sided independence assertion
requires an additional suitable soundness assumption. This report
proves neither consistency nor inconsistency.

\section{Further research questions}
\begin{enumerate}
\item \textbf{Formalize the small new interface.} The most focused target
is the ordinary rank lemma, signed auxiliary target recovery, and the
29-integer substitution theorem. They can be proved parametrically
before formalizing any astronomical fixed coefficient.
\item \textbf{Close the compiler formalization gap.} Connect an executable
finite-window coefficient compiler and its carry/population proofs to
the existing K215 endpoint. This would make the strongest current
unformalized dependency explicit in a proof assistant.
\item \textbf{Reduce the seven remaining positive coordinates.} They are
$J,F,\alpha,z,s,\eta,\zeta$. A valid reduction needs a proof that
zeros outside the positive domain are harmless, or a completeness
construction meeting a narrower polynomial image. Congruence
observations alone are not enough.
\item \textbf{Improve the product finalizer.} The degree-99 product costs
two more operations than the existing degree-155 product. Can the
ordinary strong equality be combined into one unit factor with lower
cost, without admitting a negative-unit branch?
\item \textbf{Lower the quartic variable count.} The current 33-variable
lift is explicit and verified. Eliminating a register must preserve
both quadratic degree and positivity of the inverse lift; removing
a name while creating cubic products is not an improvement of this
metric.
\item \textbf{Reduce coefficient height.} A direct compiler for the C222
calculus, a smaller local alphabet, or a different mask layout might
matter much more than one arithmetic gate when coefficient bits are
charged. A 214-state TM alone would not necessarily improve the
arithmetic skeleton.
\item \textbf{Compare with explicit low-variable universal forms.}
Instantiating a classical nine-variable or modern integer universal
construction with a fully bound proof predicate would provide a useful
alternative. Its degree, coefficient recipe, arithmetic circuit, and
formalization assumptions should all be disclosed together.
\end{enumerate}

\appendix
\section{A closed specification of the four coefficients}\label{app:coefficients}
This appendix retains package 06's fixed recipe. The accompanying
\src{vendor/blank67/compiler_constants.py},
\src{vendor/blank67/blank_helical.py}, and \src{data/k215.tm}
are exact source copies with hashes. No coefficient from another
package is substituted.

\subsection{The local alphabet and the two distinguished windows}
K215 has working states $0,\ldots,214$, halt state 215, and two symbols.
For tableau padding only, make state 215 absorbing without changing
the transition that enters it. Ordinary cell codes are
\[
 2t+b\quad(b\in\{0,1\},\ 0\le t\le216),
\]
where $t=0$ means headless and $t>0$ means state $t-1$.
Add initial phases $L,Q,R$ with codes 434,435,436, and boundaries
$V_b,H_b$ with codes 437,438. Their phase payloads are respectively
blank/headless, blank/state 0, and blank/headless.
Thus the alphabet size is $a_*=439$.

Every $3\times3$ window is read in row-major order at time offsets
$-1,0,+1$. A center-valid window satisfies these local clauses:
\begin{enumerate}
\item Vertical-boundary bits agree between center and following-time
cell. No two $V_b$ are horizontally adjacent and no two $H_b$ are
vertically adjacent.
\item If neither center nor right neighbor is $V_b$,
their horizontal-boundary bits agree.
\item An interior cell has an initial phase exactly when its predecessor
is $H_b$. Adjacent phases are $LL,LQ,QR,RR$; a first interior cell
next to $V_b$ is $L$, and a last one is $R$.
\item A head immediately before a following $H_b$ is in state 215.
No head moves out through a side boundary.
\item If the following cell is interior, its payload is the exact
deterministic three-neighbor K215 update. Multiple arriving heads are
forbidden. Boundary neighbors act as blank headless cells.
\end{enumerate}
The executable predicate fixes these clauses for boundary centers as
well; its exact Python source is part of the coefficient definition.
The first window is
\[
 (438,438,438,434,435,436,0,1,188),
\]
the actual first transition around the initial head. The second is
\[
 (0,438,0,0,437,438,0,437,0).
\]
It is center-valid but globally impossible: its predecessor has
vertical bit 0 and its center vertical bit 1, violating the predicate
at that predecessor. Enumerate these two first, then every other
center-valid window lexicographically. The exact cardinality is
\[
 k_*=1\,420\,903\,358\,357\,108\,792\,377.
\]
It equals $38\,256\,607\,897\cdot439^4$; the four corners are
unconstrained. The included cardinality checker verifies the five-cell
cross count against the local predicate by a finite partition.

\subsection{Numerical layout}
The following fixed integers specify the layout completely:
\begin{center}\small
\begin{tabular}{lr}
\toprule
Symbol & Value\\
\midrule
$\ell$ & $71$\\
$j_{\max}$ & $710\,451\,679\,178\,554\,395\,461$\\
$e_d=k_*+12a_*$ & $1\,420\,903\,358\,357\,108\,797\,645$\\
$M=k_*+15a_*+1$ & $1\,420\,903\,358\,357\,108\,798\,963$\\
$E_{\max}=27M$ & $38\,364\,390\,675\,641\,937\,572\,001$\\
$H_*$ & $72\,466\,071\,276\,212\,548\,748\,431$\\
$T_1$ & $149\,194\,852\,627\,496\,423\,892\,873$\\
$T_2$ & $302\,652\,415\,330\,064\,174\,180\,879$\\
$b$ & $5^{33}$\\
$L_*$ & $5^{34}$\\
$d=bL_*$ & $5^{67}$\\
\bottomrule
\end{tabular}
\end{center}
Define $A_*=2^{71}$, $r_*=2^b$, $B=r_*^{L_*}=2^d$, and
\[
 \mu=A_*-2+\sum_{j=1}^{j_{\max}}A_*^j.
\]
For $u,v\in\{0,1,2\}$ and $0\le t<a_*$ put
\[
 e(u,v,t)=k_*+3a_*+(2-u)3a_*+(2-v)a_*+t,
 \qquad j(u,v,t)=1+(3u+v)a_*+t.
\]
Let $\mathcal P$ contain the selector positions $0,\ldots,k_*-1$,
all positions $e(u,v,t)$, the single dummy position $e_d$, and the
four anchors $M,3M,9M,27M$. These positions are distinct.
For the $i$th window $w_i$ define
\[
 c_i=1+\sum_{u,v=0}^2 A_*^{j(u,v,w_i[3u+v])}
       +\sum_{t=1}^4 A_*^{9a_*+t}.
\]
For payload positions put $c_{e(u,v,t)}=A_*^{j(u,v,t)}$;
for the four anchors in their displayed order put
$c_{e_t}=A_*^{9a_*+t}$; put $c_{e_d}=0$.
Now set
\begin{align*}
 D_C={}&r_*^{3a_*}+r_*^{H_*+a_*}+r_*^{8M}+r_*^{24M}
  +\sum_{e\in\mathcal P}c_e(r_*^{T_1-e}+r_*^{T_2-e}),\\
 D_R={}&r_*^{H_*},\\
 M_C={}&B-1-\sum_{e\in\mathcal P\setminus\{1\}}r_*^e-2r_*^{e_d},\\
 M_F={}&\mu(r_*^{T_1}+r_*^{T_2})
 +\sum_{u=0}^{1}\sum_{v=0}^{2}\sum_{t=0}^{a_*-1}r_*^{e(u,v,t)}\\
 &+\sum_{u=0}^{2}\sum_{v=0}^{1}\sum_{t=0}^{a_*-1}r_*^{H_*+e(u,v,t)}
 +r_*^{9M}+r_*^{27M}+4.
\end{align*}
The four coefficients used everywhere in this report are exactly
\[
 \boxed{\beta=B-1,\quad K=D_C+BD_R,\quad M_C,\quad M_F^*=M_F+B-1.}
\]
All summation ranges and exponents here are fixed finite integers.
The source table, local predicate, ordering, and formulas specify
unique integers even though expanding them is infeasible.

\subsection{Why this recipe encodes a blank halting computation}
The two clause bands impose local selector/payload consistency and
overlap; the four anchors align the temporal shift with whole cells.
The mask population is $d$, $M_C\equiv2\pmod4$, and
$M_F\equiv4\pmod8$. The chosen spacing makes
$2D_C-D_R\equiv3\pmod5$, so dummy-bit tuning is possible without
a coefficient-dependent correction bit.

A marked periodic valid tableau contains a finite strip between
vertical boundaries. The initial phase graph has no cycle through
$Q$, so boundaries occur on both sides of Start. Periodicity gives
a later horizontal boundary, and the rectangle between successive
boundaries follows K215 from one blank initial head to halt.
Conversely a finite halting computation can be padded independently
in width and height, including to powers of five. Helical seams
cross time-independent vertical boundaries, and preserve every local
clause. This supplies the canonical $N$ used in the square-scale
argument. Full inherited proofs, not just this synopsis, are included
in \src{dependencies/blank_helical_proof.md} and the compiler notes.

\section{Reproduction and artifact map}
From the extracted archive directory run
\begin{lstlisting}
python3 src/reproduce.py
pdflatex -interaction=nonstopmode -halt-on-error report.tex
pdflatex -interaction=nonstopmode -halt-on-error report.tex
\end{lstlisting}
The reproduction driver checks source hashes, rebuilds the arithmetic
packets in a temporary directory, compares them with the delivered
packets, runs the independent circuit and quadratic audits, and runs
the selected finite Pell and compiler checks. It does not download
dependencies, build proof assistants, or expand giant coefficients.
\begin{center}\small
\begin{tabular}{p{53mm}p{93mm}}
\toprule
Path & Contents\\
\midrule
\src{report.tex}, \src{report.pdf} & Complete manuscript and compiled PDF.\\
\src{src/} & Standard-library generators and independent verification driver.\\
\src{data/} & Literal K215 table; emitted positive, integer, and quartic specifications.\\
\src{proofs/} & Ordinary-root proof, signed-domain proof, independent sign review, exact integer degrees.\\
\src{vendor/blank67/} & Exact coefficient recipe, local predicate, and parent arithmetic packets.\\
\src{dependencies/} & Pinned ordinary mathematical proof dependencies and comparison record.\\
\src{SOURCE_MANIFEST.json} & Source paths, snapshot identities, and hashes.\\
\src{VERIFICATION.json} & Results and scope of checks performed for this delivery.\\
\src{licenses/} & Inherited repository and machine-source notices.\\
\bottomrule
\end{tabular}
\end{center}
The deliverable is prepared for repository incorporation; no remote
repository change is part of this archive delivery.

\begin{thebibliography}{99}\small
\bibitem{reconciliation} ProveIt, \emph{Con(ZFC) in Diophantine form},
eight-package reconciliation, snapshot \pin.
\href{https://github.com/VladimirReshetnikov/ProveIt/tree/3b201cdea32d2c816f5f3d20734051f25fd3ef14/SetTheory/Cardinals/docs/reports/hilbert-tenth-problem/con-zfc-diophantine-forms}{Pinned repository collection}.
\bibitem{blank67} ProveIt, package 06, \emph{Explicit Diophantine equations
for Con(ZFC)}, \src{report.tex}, and its proof notes, 22 October 2026.
Located in the collection of reference~\cite{reconciliation}.
\bibitem{ordinary02} ProveIt, package 02, \emph{A fixed Diophantine
representation of Con(ZFC)}, \src{report.tex} and
\src{PELL_INTERFACE_LEMMAS.md}, 22 October 2026.
\bibitem{start08} ProveIt, package 08, \emph{Explicit Diophantine equations
for the consistency of ZFC}, and
\src{start_only_quartic34_proof.md}, 22 October 2026.
\bibitem{rank} ProveIt, \emph{Recovering integrality after relaxing an
auxiliary Pell norm},
\src{Computability/HilbertTenthProblem/Papers/1980/PELL_RELAXED_AUXILIARY_PROOF.md}.
Copied unchanged in this archive's \src{dependencies/}.
\bibitem{stepdown} ProveIt,
\src{EXPLORATION_FIXED_MINUS_INDEX_PARITY.md},
\src{HALF_PARAMETER_PELL_92_PROOF.md}, and
\src{complete75_reversed_auxiliary89.md}.
Pinned source copies are in \src{dependencies/}.
\bibitem{kernel} ProveIt,
\src{pell_kernel_half_binomial42.md},
\src{complete75_asymmetric_scale_tradeoffs.md}, and
\src{complete75_half_binomial_compiler.md}, under
\src{Computability/HilbertTenthProblem/Papers/research-wip/native-stream-queue/}.
\bibitem{k215} ProveIt, \emph{Formal correctness of the exact 215-state machine},
\href{https://github.com/VladimirReshetnikov/ProveIt/tree/3b201cdea32d2c816f5f3d20734051f25fd3ef14/Computability/BusyBeaver/Research/ZFC/k225-formal/verification/k215}{Pinned endpoint verification directory}.
\bibitem{jones} J. P. Jones, \emph{Universal Diophantine equation},
Journal of Symbolic Logic \textbf{47} (1982), 549--571,
doi:10.2307/2273588. The repository includes a corrected annotated transcription.
\bibitem{bayer} J. Bayer, M. David, M. Ha\ss ler, Y. Matiyasevich,
and D. Schleicher, \emph{Diophantine equations over $\mathbb Z$:
universal bounds and parallel formalization}, arXiv:2506.20909v2 (2025).
\url{https://arxiv.org/abs/2506.20909}.
\end{thebibliography}
\end{document}
